\scanpage{17}
1602 --- Série A \hfill C. R. Acad. Sc. Paris, t. 278 (24 juin 1974)

\figure{\begin{tikzpicture}\draw (0,0)--(12,0);\end{tikzpicture}}

$2^\circ\ \mathrm{A}=\mathbf{Z}[i]$ :
\[
\operatorname{Br}(\mathbf{Z}[i])=\mathrm{Q}(\mathbf{Z}[i])=0\qquad\text{et}\qquad r=0.
\]
D'où $\mathrm{HW}(\mathbf{Z}[i])=\operatorname{Witt}(\mathbf{Z}[i])=0$.

(*) Séance du 4 février 1974.

(${}^1$) M. \textsc{Auslander} et O. \textsc{Goldmann}, \textit{Trans. Amer. Math. Soc.}, 97, 1960, p. 376-409.

(${}^2$) H. \textsc{Bass}, \textit{Topics in algebraic K-theory}, Tata Institute of Fundamental Research. Bombay, 1966.

(${}^3$) A. \textsc{Fröhlich}, \textit{Quart. J. Math.}, Oxford, (2), 22, 1971, p. 401-423.

(${}^4$) A. \textsc{Grothendieck}, \textit{Le groupe de Brauer III. Dix exposés sur la cohomologie des schémas}, North-Holland, 1968.

(${}^5$) D. \textsc{Harrison}, \textit{Abelian Extensions of Commutative Rings}, A. M. S. Memoirs, 1965, p. 68-79.

(${}^6$) O. \textsc{Laborde}, \textit{Bull. Sc. math.}, $2^{\mathrm e}$ série, 96, 1972, p. 199 à 208.

(${}^7$) A. \textsc{Larotonda}, A. \textsc{Micali} et O. E. \textsc{Villamayor}, \textit{Sur le groupe de Witt}, Istituto Nazionale di Alta Matematica (\textit{Symposia Mathematica}, XI, 1973, p. 211 à 219).

(${}^8$) J. \textsc{Milnor}, \textit{On simply connected manifolds} (Symposium de Mexico, 1958, p. 122-128).

(${}^9$) Ph. \textsc{Revoy} (à paraître).

(${}^{10}$) J. P. \textsc{Serre}, \textit{Corps locaux}, Hermann, Paris, 1962.

(${}^{11}$) C. \textsc{Small}, \textit{Trans. Amer. Soc.}, 156, 1971, p. 455-491.

\begin{flushright}
\textit{U. E. R. de Mathématiques,\\
Université des Sciences et Techniques\\
du Languedoc,\\
place Eugène-Bataillon,\\
34060 Montpellier.}
\end{flushright}
\mspage{16}
